\documentclass[10pt,a4paper]{amsart}
\usepackage[margin=19mm]{geometry}
\usepackage{amsmath,amssymb,mathtools}
\usepackage[scr=rsfs,cal=euler]{mathalfa}
\usepackage{xfrac}
\usepackage[unicode,hidelinks,bookmarks=false]{hyperref}
\newcommand{\ie}{i.e.\ }
\newcommand{\cf}{cf.\ }
\newcommand{\resp}{respectively\ }
\newcommand{\Subsection}{\subsection}
\providecommand{\nobreakdash}{\nobreak\hbox{-}}
\setlength{\emergencystretch}{2em}
\multlinegap0pt
\usepackage{bm}
\usepackage{bbm}
\usepackage{dsfont}
\usepackage{xspace}
\usepackage{cancel}
\usepackage{mathtools}
\usepackage{amsthm}
\makeatletter
\@ifundefined{definition}{\newtheorem{definition}{Definition}}{}
\@ifundefined{lemma}{\newtheorem{lemma}[definition]{Lemma}}{}
\@ifundefined{prop}{\newtheorem{prop}[definition]{Proposition}}{}
\makeatother
\makeatletter
\@ifundefined{dmc}{\newenvironment{dmc}[1]{\renewcommand{\Title}{#1}\begin{defititle}} {\end{defititle}}}{}
\makeatother
\title[Equivariant rigidity of complex and quaternionic moment--ang]{Equivariant rigidity of complex and quaternionic moment--angle manifolds}
\author{Ioannis Gkeneralis}
\date{Source: arXiv:2604.18455v2 (CC BY 4.0)}
\begin{document}
\maketitle

\noindent\emph{Source and licence.} Equivariant rigidity of complex and quaternionic moment--angle manifolds. Version arXiv:2604.18455v2 (CC BY 4.0), distributed under the Creative Commons Attribution 4.0 International licence, as recorded in the arXiv licence metadata for this exact version and shown on the abstract page https://arxiv.org/abs/2604.18455v2. Full rights notice: licenses/arxiv-2604.18455-CC-BY-4.0.txt.\par\smallskip

\noindent\textit{Editorial scope.} Counted unit: the Prop whose source label is \texttt{prop:quotient\_manifold} in the article (source0.tex lines 1412--1425, source label \texttt{prop:quotient\_manifold}), delivered together with the article's own complete original proof of it (source0.tex lines 1427--1476, one \texttt{proof} environment of 50 lines). No mathematics is written, repaired or re-derived: statement and proof are the article's own text line for line. Required context retained uncounted: lem:Kfree (lines 1349--1351); lem:dim\_match (lines 1388--1393). Every definition, display and label of those lines is retained unchanged. Omitted from this package and not redistributed: 1--1348, 1352--1387, 1394--1411, 1426--1426, 1477--1851. Nothing inside a retained excerpt is omitted or altered: every retained body is delivered exactly as the article's lines are written, so no omission receipt of any kind accompanies this package. Citations: the retained excerpts carry no citation command at all, so no bibliography entry of the article is transcribed and no reference is redistributed. Reference adaptation: none was needed and none was made; every reference inside the delivered unit resolves inside it, so no unresolved reference or citation is produced. Printed slips kept literal: no printed word, symbol, hypothesis or number of the counted statement or of its proof was altered. Layout and class: the article's own class is \texttt{article} with packages including etex, inputenc, babel, lmodern, textcomp, verbatim, appendix, cite, amsmath, amsthm, amsfonts, amssymb, mathtools, stmaryrd; the wrapper is the lane's pinned \texttt{amsart} editorial wrapper at 10pt on A4, which restates the theorem-like environments the retained text needs and adds the article's own macro definitions that the retained text needs (none), with extra wrapper packages (bm, bbm, dsfont, xspace, cancel, mathtools). The delivered numbering is the unit's own. Assets: no retained span contains a figure environment, a graphics-inclusion command, a PDF-page inclusion, a table or a TikZ picture; no image file is copied into this package, the package declares no asset dependency and no image file is redistributed with it. Licence: version arXiv:2604.18455v2 of this article is distributed under the Creative Commons Attribution 4.0 International licence, as recorded in the arXiv licence metadata for this exact version and shown on the abstract page https://arxiv.org/abs/2604.18455v2. A full rights notice is delivered at licenses/arxiv-2604.18455-CC-BY-4.0.txt. Characters: every retained excerpt is pure ASCII, so the delivered engine, XeLaTeX with its default Latin Modern text font, reports no missing character.
\begin{lemma}\label{lem:Kfree}
The subgroup $K$ acts freely on $N$.
\end{lemma}

\begin{lemma}\label{lem:dim_match}
Let \(X\) and \(Y\) be closed topological manifolds, and suppose that \(f:X\to Y\) is a homotopy equivalence. Then
\[
    \dim X=\dim Y.
\]
\end{lemma}

\begin{prop}\label{prop:quotient_manifold}
Let \(G=T^m\) in the complex case and \(G=Q^m\) in the quaternionic case.  Let \(N\) be a closed locally linear \(G\)-manifold equipped with a \(G\)-equivariant homotopy equivalence
\[
    f:N\longrightarrow \mathcal Z_P
\]
in the complex case, or
\[
    f:N\longrightarrow \mathcal Z_P^{\mathbb H}
\]
in the quaternionic case. Let \(K\subset G\) be the kernel subgroup associated to the characteristic data. By Lemma~\ref{lem:Kfree}, \(K\) acts freely on \(N\). Then \(N/K\) is a closed topological manifold and the residual \(G/K\)-action on \(N/K\) is locally linear.

Moreover, \(\dim(N/K)=2n\) in the complex case, and
\(\dim(N/K)=4n\) in the quaternionic case.
\end{prop}

\begin{proof}
Since \(K\) is a compact Lie group acting freely and locally linearly on the closed manifold \(N\), the quotient map
\[
    N\longrightarrow N/K
\]
is a locally trivial principal \(K\)-bundle. In particular, \(N/K\) is a closed topological manifold and
\[
    \dim(N/K)=\dim N-\dim K.
\]

By Lemma~\ref{lem:dim_match}, in the complex case
\[
    \dim N=\dim\mathcal Z_P=m+n.
\]
Since \(K\cong T^{m-n}\), we obtain
\[
    \dim(N/K)
    =
    (m+n)-(m-n)
    =
    2n.
\]

In the quaternionic case,
\[
    \dim N=\dim\mathcal Z_P^{\mathbb H}=3m+n.
\]
Since \(K\cong Q^{m-n}=(S^3)^{m-n}\), we have
\[
\dim K=3(m-n),
\]
and therefore
\[
\dim(N/K)=(3m+n)-3(m-n)=4n.
\]

It remains to check local linearity of the residual action. Let \([x]\in N/K\), and let \(G_x\) be the isotropy subgroup of \(x\) in \(G\). Since \(K\) acts freely on \(N\), we have
\[
G_x\cap K=\{1\}.
\]
By local linearity of the \(G\)-action on \(N\), there is a \(G\)-invariant neighbourhood of the orbit \(Gx\) of the form
\[
    G\times_{G_x} V,
\]
where \(V\) is a linear \(G_x\)-representation.  Quotienting by the free \(K\)-action gives
\[
(G\times_{G_x}V)/K\cong (G/K)\times_{G_x}V,
\]
where \(G_x\) is identified with its image in \(G/K\). This is a linear tube for the residual \(G/K\)-action near \([x]\). Hence the residual action of \(G/K\) on \(N/K\) is locally linear.
\end{proof}

\end{document}
